\documentclass[11pt,a4paper]{article}
\usepackage[T1]{fontenc}
\usepackage{lmodern}
\usepackage[margin=2.5cm]{geometry}
\usepackage{amsmath}
\usepackage{graphicx}
\usepackage{booktabs}
\usepackage{authblk}
\usepackage[round]{natbib}
\usepackage[hidelinks]{hyperref}

\title{Title of the Review: The State of a Research Field}
\author[1]{First Author}
\author[2]{Second Author}
\affil[1]{Department, University, City, Country}
\affil[2]{Institute, City, Country}
\date{}

\begin{document}

\maketitle

\begin{abstract}
Say what the review covers, why now, the main themes and the open questions in
a few sentences.
\end{abstract}

\noindent\textbf{Keywords:} first keyword, second keyword, third keyword

\section{Introduction}
Define the field and the scope of this review, and how the literature was
selected. Earlier reviews include \citet{lecun2015deep}.

\section{Background}
Introduce the concepts the rest of the review relies on~\citep{shannon1948}.

\section{First theme}
Summarise and compare the main lines of work. A table helps readers compare them.

\begin{table}[htbp]
  \centering
  \caption{Approaches compared in this review.}
  \label{tab:approaches}
  \begin{tabular}{lrr}
    \toprule
    Approach & Since & Data needed \\
    \midrule
    Approach A & 2015 & Large \\
    Approach B & 2021 & Small \\
    \bottomrule
  \end{tabular}
\end{table}

\section{Second theme}
Continue with the next group of studies~\citep{he2016deep,vaswani2017attention}.

\section{Open questions and outlook}
State what remains unknown and where the field is heading.

\section*{Acknowledgements}
Thank the people and funding that supported this work.

\bibliographystyle{plainnat}
\bibliography{references}

\end{document}
